\section{Threats to Validity}
\label{sec:threats}

The threats taxonomy follows Wohlin et al.~\cite{wohlin2012experimentation}.

\textbf{Conclusion validity.} The evaluation comprises 2{,}250
P@50 observations summarised with 95\% percentile bootstrap
CIs (10{,}000 resamples). The H4 margin ($+0.7$~pp) is within
the bootstrap noise at most cell-windows; we report it as a
descriptive cell-mean win rather than a sharp claim. H3's
Spearman $\rho$ is computed on three points (one per $K$) and
is statistically under-powered; the rejection is treated as a
descriptive negative rather than a strong inference.

\textbf{Internal validity.}

\textit{(i) Selection-policy coupling.} As in Papers~7 and~8, all
five strategies share the fleet trajectory generated by
HygienePrio-full under \textsf{fixed} weights. The strategies
differ only in scoring weights, isolating the recalibration
decision. Paper~9~\cite{paper9selftraj} addressed the
selection-policy coupling threat separately for the cross-method
comparison; the present paper inherits that finding (selection
matters) and bounds claims to the fixed-policy trajectory.

\textit{(ii) Threshold $\tau = 0.05$ is fixed.} The pre-registration
locked $\tau$ to avoid post-hoc tuning. The reported H1 success and
H2 cost are specific to this threshold; a $\tau$ sweep is future work.

\textit{(iii) Detector simplicity.} The one-threshold magnitude test
is the minimum-complexity detector. CUSUM and Bayesian
alternatives might produce different fire patterns and shift the
H1/H2 balance.

\textbf{Construct validity.} P@50 at $K = 200$ becomes noisy at
late windows due to small positive-set sizes; the H4 margin is
small and the cell-mean comparison is the appropriate aggregation.

\textbf{External validity.} Inherited from Papers~4--8: synthetic
EEHDA structural priors, fixed $\lambda = 3$, 14-day window, 0.15
patch debt reduction. Real fleets with different rates of
per-window distributional shift would shift the H1/H2 trade-off
curve; the qualitative result (a simple detector solves H1 but
pays H2) is more robust than the specific cell-mean numbers.
